\chapter{La problématique de financement de l'open source}

Beacoup de projets libres sont gratuits, mais cela ne signifie pas forcément
que rien n'a été investi : il existe de nombreux modèles de financement des
projets opensource dont certains initiés par des sociétés comme RedHat ou Apple
qui sont en croissance constante.\\ Le logiciel libre peut donc être rentable à
condition de savoir exploiter le modèle de financement adapté à l'ampleur du
projet.

\section{Quelques modèles alternatifs}

    \subsection{Les dons}

Le financement de projets par dons est très répandu.  Si ce modèle peut servir
d'appoint aux développeurs, il n'est viable que dans très peu de cas.
    
    \subsection{Donations par micropaiement}

Le modèle du pay per like est une version plus organisée du modèle par dons.
Les dons se font non plus directement au projet mais passent par un oganisme
chargé de répartir équitablement un budget mensuel, défini par le donneur,
entre tous les projets qu'il aura explicitement marqués dans le mois.
\url{flattr.com}, par exemple, permet distribuer le budget entre les projets
marqués sur \url{github.com}.\\ Flattr, bien que compatible pour le financement
de projets informatiques est surtout prévu pour récompenser les créateurs de
contenu artistique.
    
    \subsection{Le Crowdfunding}

\paragraph{} En français, on l'appelle "financement collaboratif". Il s'agit
d'un model de financement basé sur la collecte de don. Le leader sur le marché
est le site Kickstarter. Il propose à un utilisateur ou un groupe d'utilisateur
de mettre en ligne l'idée de son projet pour qu'il apparaisse sur le site. Les
visiteurs peuvent alors consulter cette idée et choisir ou non d'investir de
l'argent pou r la voir se développer. Les personnes ayant contribuer au
développement sont très souvent récompensé lors de la fin du projet. Il peuvent
par exemple rencontrer les créateurs du projet qu'ils ont soutenu, obtenir des
goodies\footnote{\textbf{goodies :} produits dérivés faisant référence à un
film, un jeu-vidéo...} ou même pouvoir voter pour l'implémentation de
fonctionnalités.  Ce mode de financement développe l'entraide et donne la
capacité financière afin de mener à bien un projet prometteur.

    \subsection{Licence Double}

La licence double est un modèle adopté notamment par MySQL consiste à permettre
l'utilisation de son produit différemment suivant l'usage. Dans le cas de
MySQL, l'utilisation personnelle ou l'inclusion dans un projet sous licence GPL
est gratuite tandis qu'une utilisation profesionnelle est payante.

    \subsection{Modèle Redhat}

Redhat distribue son produit gratuitement mais propose un support payant
garantissant la stabilité de son produit. Ce modèle a été adopté par
l'entreprise en 2003, qui depuis cette date est en constante croissance.  C'est
un modèle similaire qui régit Canonical.

    \subsection{Modèle Apple}

Le système d'exploiation OSX d'apple est dérivé d'un système ouvert.\\ La
stratégie d'apple pour son système d'exploitation OSX a consisté à changer la
licence et apporter des modifications pour ensuite le vendre sur ses machines.

    \subsection{Modèle Bountysource}

\url{bountysource.com} est un site d'annonces de demande de code. Les
utilisateurs peuvent demander des fonctionnalités et se regrouper à plusieurs
sur une même annonce pour augmenter la prime. Ainsi, plus une fonction
intéresse d'utilisateurs, plus la somme proposée pour son implémentation est
élevée et donc susceptible d'intéresser un développeur.

    \subsection{Modèle Pay What You Want}
    
    \paragraph{} Le modèle Pay What You Want consiste à proposer un produit à
    un utlisateur pour la somme qu'il estime nécessaire de payer. Il peut y
    avoir un seuil minimum à verser ou bien non. Dans le deuxième cas,
    l'utilisateur peut donc acquérir le produit gratuitement de façon légale,
    mais il peut aussi verser un montant qu'il juge adapté au produit de façon
    à récompenser le ou les concepteurs pour leur travail.

    \paragraph{} Ce type de modèle s'est démocratisé en 2007 grâce au groupe de
    musique Radiohead qui a proposé sur son site officiel son nouvel album
    intitulé "in rainbow" afin de lutter contre le piratage. On remarque
    l'utilisation de ce modèle de financement et de ses nombreux dérivés (Pay
    what you wish, Pay what you can etc.) surtout dans le secteur des
    logiciels, de la musique, de la restauration et de l'hotellerie. On peut en
    voir un parfait exemple à l'adresse suivant
    \url{https://www.humblebundle.com/}.


\section{Gittip confronté aux autres modèles}

\subsection{La stabilité}


%A voir plus tard // Besoin de Trier%

Gittip est un modèle proche de celui de Flattr à la différence que les
donnations sont obligatoirement récurrentes et faites à un développeur ou une
équipe plutot qu'un projet précis.

\subsection{Les modèles proches}

\paragraph{} Le caracère récurrent des dons assure plus de stabilité que les
dons ponctuels via Flattr ou Paypal.

\paragraph{} Les équipes financées grace a Gittip dépendent d'un grand nombre
d'acteurs, ce qui leur permet de mesurer leur popularité et de ne pas dépendre
d'un seul investisseur qui pourrait décider d'abandonner le projet.

\paragraph{} Gittip est lui même financé via Gittip et non via le prélèvement
sur les dons fait aux développeurs, le seuls prélèvements effectués servent à
couvrir les opérations bancaires.

\begin{figure}[h!] \captionsource{Comparatif des prélèvements par
    plateforme}{\url{https://www.gittip.com/about/charts.html}} \begin{center}
        \begin{tabular}{|c|c|c|} \hline % Permet d'avoir une séparation via un
            trait \textbf{Paypal} & \textbf{Gittip} & \textbf{Flattr} \\ \hline
            2,2\% + 30\textcent & 2,9\% + 30\textcent & 10\% \\ \hline
        \end{tabular} \end{center} \end{figure}

Dans le tableau ci-dessus, le prix des transactions pour utilisé Paypal est
celui appliqué au bouton "Donate".


\paragraph{} Gittip permet également de gérer efficacement son budget à
l'interieur d'une équipe grâce au mécanisme précédemment décrit, ce que ne
gèrent pas flattr ni Paypal.

\subsection{Les autres modèles}

Les modèles restants visent soit à lancer un projet, soit à financer des
projets de grande ampleur, et donc nécessitant un investissement plus
important.  Actuellement, Chad Whitacre gagne 23300\$ par an alors que le
salaire moyen de développeur est de 90000\$ par an. Dans l'état actuel des
choses, Gittip n'est donc pas viable en tant que source principale de revenus.
Toutefois, il est important de noter que la somme échangée par semaine augmente
avec le nombre d'utilisateurs actifs depuis la création du site.

\begin{figure}[h!] \captionsource{Nombre d'utilisateurs en fonction du
temps}{\url{https://www.gittip.com/about/charts.html}} \centering
\includegraphics[width=16cm]{images/gittip_amount.eps} \end{figure}
\begin{figure}[h!] \captionsource{Somme échangée au sein de Gittip (en \$) en
    fonction du temps}{\url{https://www.gittip.com/about/charts.html}}
\centering \includegraphics[width=16cm]{images/gittip_user.eps} \end{figure}

\paragraph{} En effet, la différence essentielle entre des sites de financement
tels que KickStarter et Gittip consiste à l'utilisation des fonds récoltés. Du
côté de KickStarter, l'argent est versé par des personnes soutenant un projet
et donc cet argent ne pourra être utilisé que pour mener à bien la création de
ce projet. En revanche, pour Gittip, l'argent récolté ressemble davantage à un
salaire que la personne recevra toute les semaines de façon prévisible et cet
argent pourra être utilisé par la personne comme bon lui semble. L'intêret
étant bien entendu de fournir une aide financière à la personne pour qu'elle
puisse consacrer un maximum de temps à ses projets personnels.

\paragraph{} parler de : pas de financement important (comparé a kickstarter
pour initier un projet ou le financement par une société tierce comme chez
mozilla)


